\chapter{Theoretical Analysis}
This chapter introduces the foundational concepts necessary to understand the proposed data cleaning framework. 
It covers the problem statement, definitions of relational data, types of errors, the fundamentals of LLMs, evaluation metrics and the supporting roles of knowledge graphs and data profiling. 
\section{Problem Statement}
Errors in datasets significantly degrade data quality and undermine downstream tasks such as machine learning, statistical analysis and decision-making processes. 
This thesis focuses on the problem of automated data cleaning in tabular datasets by detecting and repairing erroneous cells.\newline
\indent Let a tabular dataset be defined as \( D = \{r_1, r_2, \dots, r_M\} \), where each \( r_i \) represents a row consisting of \( N \) attributes \( [a_1, a_2, \dots, a_N] \). 
Each cell in the dataset can be denoted as \( D_{ij} \), corresponding to row \( i \) and column \( j \).\newline
\indent Given a dirty dataset \( D^{d} \), containing unknown errors and an unknown ground truth \( D^{t} \), an error is formally defined as any cell \( D^{d}_{ij} \) for which \( D^{d}_{ij} \ne D^{t}_{ij} \). 
The objective of this thesis is to develop a framework that automatically detects such errors and repairs them with corrected values \( D^{*}_{ij} \) such that \( D^{*}_{ij} \approx D^{t}_{ij} \), with the ideal case being \( D^{*}_{ij} = D^{t}_{ij} \).
\section{Concepts}
\subsection{Relational Datasets}
A relational or tabular dataset is a structured collection of records, where each record is represented as a row in a table. 
Each row is a tuple of attribute values aligned with a fixed schema. 
For example, a customer dataset may have the schema [CustomerID, Name, Age, Country], with a data instance such as \{'12345', 'Jane Doe', 55, 'USA'\}.\newline
Relational datasets are a preferred data format in enterprise systems due to their structured nature, ease of access and support for query languages such as SQL \cite{SCRD}. 
Unlike non-relational datasets, which are more flexible in structure but often less semantically rich, relational datasets enable clear definitions of schema, constraints and inter-column relationships. 
These properties are essential for automated data cleaning methods.

\subsection{Error Types in Tabular Datasets}
Data quality problems in tabular datasets often appear as various types of errors. 
In this thesis, we focus on the most common and important error types found in real-world data:
\newline
\textbf{(Disguised) Missing Values:} These include both standard missing entries (e.g., NULL, NaN) and values that act as placeholders (e.g., 'X', -999, 'Unknown'). 
Detecting standard missing values is trivial, while DMVs are more challenging due to their arbitrary formats and context dependence. 
\newline
\textbf{Outliers:} These are values that significantly deviate from the rest of the data. 
For instance, a salary that is ten times higher than the average might indicate an anomaly. 
\newline
\textbf{Formatting Errors:} These errors occur when values do not match the expected representation.
This category includes pattern violations, such as values that deviate from a dominant format (e.g., an ID written as '123-456-789' instead of the expected '1234-5678'), incorrect data types, such as strings in numeric columns, and duplicate or concatenated values (e.g., 'Smith, Smith').
\newline
\textbf{Functional Dependency Violations:} An FD of the form A → B implies that column A should uniquely determine B. 
Violations occur when this constraint is broken, such as when a single zipcode maps to multiple cities (e.g., 'Utrecht' → 030, 'Amsterdam' → 030, while City → ZipCode should hold).
\newline
\textbf{Semantic Errors:} These values are syntactically correct but contextually incorrect, such as typos, unrealistic entries (e.g., a date of 45/45/2900 or an age of 130) or inconsistent naming (e.g., 'Main Street' vs 'Main St.'). 
Such errors are hard to detect using traditional statistical methods and often require semantic understanding.
\newline\newline
It is important to note that for certain error types, particularly outliers and functional dependency violations, traditional data cleaning approaches often rely on naive, uniform strategies, such as removing all outliers or enforcing a single repair for all violations.  
These approaches typically ignore valid exceptions that arise from contextual or real-world semantics. 
As a result, correct values may be incorrectly modified or removed.
Incorporating semantic reasoning and contextual awareness helps distinguish true errors from legitimate exceptions, reducing overcorrection and preserving valid data.

\subsection{Context for Data Cleaning}
Effective error detection and repair depend heavily on contextual understanding. 
Context in tabular datasets can be considered at multiple levels, each providing different signals for identifying and correcting errors:
\newline
\textbf{Value-level context:} Information about individual cell values, such as their datatype, format, or how unusual a value is relative to the expected column type. 
For example, a string value in a numeric column or a negative value in a salary column may indicate an error.
\newline
\textbf{Row-level context:} Captures relationships between values within the same row and helps identify inconsistencies across related columns.
For example, a row with Country = 'USA' and CountryCode = 31 may indicate an error, as 31 corresponds to the Netherlands. 
This inconsistency becomes clear only when considering both columns together.
\newline
\textbf{Column-level context:} Captures patterns and statistical properties within a single column, such as value distributions or dominant formats. 
For example, if most values in a date column follow the pattern \textit{dd/mm/yyyy}, but a few entries use \textit{yyyy/mm/dd}, the latter can be flagged as formatting errors. 
Similarly, a salary value of 200,000 may be flagged as an outlier if most values lie between 50,000 and 70,000. 
\newline
\textbf{Dataset-level context:} Describes relationships across multiple columns or the entire dataset, such as correlations, functional dependencies or global constraints. 
Using this context often requires processing large parts of the dataset, which can be computationally expensive for large datasets.
\newline\newline
Column-level context is often preferred in data cleaning, as columns typically have a single semantic meaning, a consistent datatype and a dominant format, enabling uniform and scalable cleaning operations. 
Relevant context can be provided to the LLM through clear prompts with instructions and examples, enabling context-aware reasoning without processing the entire dataset.

\subsection{Large Language Models}
LLMs are transformer-based neural networks trained on large amounts of textual data to learn language representations and reasoning patterns. 
They predict the next token in a sequence using techniques like self-attention to capture semantic, syntactic and contextual relationships \cite{COMPREHENSIVE}. 
This allows them to function like a human-like chatbot, capable of interpreting text and generating appropriate responses. \newline
LLMs are generalizable, meaning that they can be fine-tuned or prompted to perform a wide range of tasks. 
They are not bound to schemas or manual rules, making them adaptable to diverse data cleaning operations. 
Finally, LLMs can reason in a human-like manner and are particularly effective at identifying context-dependent errors, similar to how an expert would review a dataset.
\newline\newline
Despite these strengths, LLMs also introduce important reliability challenges. 
Their outputs can be non-deterministic, sensitive to prompt quality and prone to hallucinations \cite{AIALIGN}, meaning that generated responses or corrections cannot always be assumed to be correct. 
As a result, simply passing raw data to an LLM for cleaning is insufficient for reliable, automated data processing. 
Instead, effective use of LLMs in data cleaning requires a structured framework that explicitly addresses these limitations and enforces consistency and correctness in the produced outputs.
\newline\newline
Recent generations of LLMs are particularly well suited for code generation \cite{LLMCODE}, which enables a different approach to data cleaning. 
Rather than generating individual repaired values, LLMs can generate executable code that applies cleaning operations to entire columns. 
This approach offers two key advantages. 
First, it improves consistency, as the same transformation logic is applied uniformly across all values in a column. 
Second, it enhances scalability, since column-level transformations can be executed efficiently by avoiding repetitive cell-level inference. 
These properties make code generation a strong foundation for building reliable and scalable LLM-based data cleaning systems.

\subsection{Data Profiling}
Data profiling is the process of analyzing the content and structure of a dataset to extract statistical summaries and metadata \cite{PRDS}. 
It provides essential information such as column types, value distributions, frequency counts, uniqueness and null counts, as well as FDs and other inter-column relationships. 
This information offers an accurate view of the dataset's characteristics and helps the LLM better understand where to focus, leading to more targeted and effective data cleaning.
In addition, data profiling addresses several limitations of LLM-based cleaning.
Statistical properties such as distributions, frequency patterns, correlations and dependencies are difficult for LLMs to infer reliably from raw data alone. 
Explicitly providing this information reduces the reasoning burden on the LLM and supports more informed and accurate cleaning decisions. 
Moreover, statistical profiling is computationally efficient and scalable, making it well suited for large datasets.

\subsection{Assumptions}
The framework assumes that tabular datasets exhibit a certain degree of structure and redundancy at the column level, which can be leveraged to infer the intended data format. 
To reduce computational costs, the framework relies on sampling rather than processing entire datasets, assuming that sampled values are representative of the overall column content. 
As a result, rare or highly irregular errors may remain undetected.
Finally, framework assumes that most data quality issues can be effectively addressed using column-level transformations. 
While row-level context is incorporated in specific cases, the overall design prioritizes column-wise standardization over individual value repairs.